% !PRINTED 312
\weberpage{312}
%<<CONT
woraus zu schließen:
\weq{7}{\begin{aligned} v_\infty &= v_0 = \frac{-1+\sqrt5}{2}, \\ v_1 &= v_4 = \frac{-1-\sqrt5}{2}. \end{aligned}}
($v_0$, $v_\infty$ müssen, da in ihnen $q$ reell ist, positiv sein.)

Aus (3) ergibt sich sodann:
\weq{8}{\begin{aligned} w_0 &= 0, \quad w_1 = w_3 = -5-2i\sqrt5, \\ w_2 &= w_4 = -5+2i\sqrt5. \end{aligned}}
Dies müssen die Wurzeln von (5) sein für $j=12^3$ und daraus erhält man die gesuchte Resolvente in der Form:
\weq{9}{w(w^2+10w+45)^2 = j(\omega) - 12^3 = \gamma_3^2.}
Diese Gleichung läßt sich auch in die Form bringen:
\weq{10}{(w+3)^3(w^2+11w+64) = j(\omega) = \gamma_2^3,}
auf die man auch direkt kommt, wenn man, anstatt $\omega=i$ zu setzen,
\[ \omega = \frac{-1+i\sqrt3}{2} \]
annimmt.

Die gefundene Resolvente vereinfacht sich noch, wenn man [nach (9)]:
\[ z = \sqrt{w} = \frac{\gamma_3}{w^2+10w+45} \]
setzt:
\weq{11}{z^5+10z^3+45z = \gamma_3,}
oder wenn man [nach (10)]
\weq{12}{w^2+11w+64 = y^3, \quad y = \frac{\gamma_2}{w+3}}
setzt:
\weq{13}{y^5 - 40y^2 - 5\gamma_2 y - \gamma_2^2 = 0,}
eine Gleichung, die sich auch ergibt, wenn man von vornherein die Annahme macht:
\[ 2y = v_\infty v_0 + v_1 v_4 + v_2 v_3. \]

In der Gleichung (13) fehlen, wie man sieht, die dritte und die vierte Potenz der Unbekannten. Dieselbe Eigenschaft kommt auch, wie leicht nachzuweisen ist, der Gleichung zu, deren Wurzeln
%CONT>>
